% Generated by analysis/scripts/export_ten_cases.py
\subsection*{统一结果卡}
\begingroup\small
\noindent\begin{tabularx}{\textwidth}{@{}p{24mm}X@{}}\toprule
技术/模式 & SRAM DCIM；二进制6T；D6CIM型16项归约、16活动输出通道\\
逻辑粒度/聚合 & B\_S=128 Byte；B\_R=16 Byte。对齐128 bit普通写口，完成16 Byte；1024事务覆盖整矩阵\\
资源/相对R0 & 8个容量分片、总128条/活动16条HCA-BFA；R0的32项改16项；完整MAC槽覆盖静态读/保持\\
服务口径 & 无周期维护；普通SRAM需持续供电\\
主导因素/选择 & 512个完整MAC槽、64次行组选择；参考5 ns计算槽与5 ns普通写槽独立，静态连接不另付重读\\
关键对照 & 32项同钟资源对照；额外握手仅适用于边界改变\\
\bottomrule\end{tabularx}
\vspace{3mm}
\begin{center}\begin{tabular}{@{}lrrr@{}}\toprule
成对条件情景 & $\rho$ (MB/s) & $\tau$ (MB/s) & $\mathrm{RI}^{*}$\\\midrule
短预算 & 58.03 & 7273 & 0.00798\\
参考（推荐） & 49.81 & 3200 & 0.01556\\
长预算 & 24.9 & 1600 & 0.01556\\
\bottomrule\end{tabular}\end{center}
\noindent 参考原始占用：streaming 2570 ns；resident 5 ns。
\par\noindent 范围为有限成对条件情景极值，不是物理不确定性包络。1 MB=$10^6$ Byte；整矩阵16384 Byte=16 KiB。源定位见本例证据笔记；公共基线shared\_baseline。
\endgroup
\clearpage
